% ============ 17 紧合基变换定理 (原书 pp.114-117) ============
\chapter{紧合基变换定理}\label{sec:17}
\index{紧合态射}
\index{紧合基变换定理}
\index{可构造层}
\index{完备簇}
\index{大平展景}

\subsection*{拓扑中的紧合基变换定理}

本小节中，所有拓扑空间都假定 Hausdorff.

回忆紧空间在连续映射下的像是紧的，从而闭的.此外，拓扑空间 $X$ 紧，当且仅当对所有拓扑空间 $Y$，投影映射 $X \times Y \to Y$ 闭(Bourbaki, N., \emph{Topologie G\'en\'erale}, I, \S 10, Th. 1 的推论 1).

连续映射 $\varphi \colon X \to S$ 称为紧合的，如果它泛闭，即对所有连续映射 $T \to S$，投影映射 $X \times_{S} T \to T$ 闭.当 $S$ 局部紧时，连续映射 $\varphi \colon X \to S$ 紧合，当且仅当 $S$ 的每个紧子集的原像紧，这时 $X$ 也局部紧(同上，命题 7).

设 $\varphi \colon X \to S$ 是连续映射，$F$ 是 $X$ 上的层.对 $s \in S$，
\[
(R^{r}\varphi_{*}F)_s = \varinjlim_{V} H^{r}(\varphi^{-1}(V), F),
\]
其中极限遍历 $s$ 在 $S$ 中的开邻域 $V$(比较 12.2).若 $\varphi$ 闭，则对 $\varphi^{-1}(s)$ 的任何开邻域 $U$，$\varphi(X \setminus U)$ 闭，且 $V \overset{\mathrm{def}}{=} X \setminus \varphi(X \setminus U)$ 是 $s$ 的开邻域.由于 $\varphi^{-1}(V) \subset U$，看到
\[
(R^{r}\varphi_{*}F)_s = \varinjlim_{U} H^{r}(U, F),
\]
其中极限现在遍历 $\varphi^{-1}(s)$ 的所有开邻域 $U$.

\begin{lemma}\label{lem:17.1}
设 $Z$ 是局部紧空间 $X$ 的紧子集.则对 $X$ 上的任何层 $F$，有典范同构
\[
\varinjlim H^{r}(U, F) \cong H^{r}(Z, i^{*}F)
\]
($Z$ 在 $X$ 中的开邻域 $U$ 上的极限).
\end{lemma}

\begin{proof}
由于两个群在内射层上对 $r > 0$ 都为零，且把层的短正合序列变成长的正合序列，只需(由导出函子的唯一性)对 $r = 0$ 验证论断.这里使用 $Z$ 与 $X$ 的假设(见 Iversen 1986, III 2.2($r = 0$ 的情形)与 III 6.3(对所有 $r$ 的推广)).
\end{proof}

\begin{theorem}\label{thm:17.2}
设 $\varphi \colon X \to S$ 是紧合映射，其中 $S$ 是局部紧拓扑空间.对 $X$ 上的任何层 $F$ 与 $s \in S$，$(R^{r}\varphi_{*}F)_s \cong H^{r}(X_s, F)$，其中 $X_s = \varphi^{-1}(s)$.
\end{theorem}

\begin{proof}
把上面的论断结合起来.
\end{proof}

设 $\varphi \colon X \to S$ 与 $F$ 如定理中所示，$f \colon T \to S$ 是连续映射.构成纤维积图
\[
\begin{tikzcd}
X \arrow[d, "\varphi"'] & X' = X \times_{S} T \arrow[l, "f'"'] \arrow[d, "\varphi'"] \\
S & T \arrow[l, "f"']
\end{tikzcd}
\]
设 $t \in T$ 映到 $s \in S$.则 $f'$ 在纤维上定义同构 $X'_t \to X_s$，因此定理蕴含
\[
(R^{r}\varphi_{*}F)_s \cong H^{r}(X_s, F) \cong H^{r}(X'_t, f'^{*}F) \cong (R^{r}\varphi'_*(f'^{*}F))_t.
\]

\begin{theorem}\label{thm:17.3}
在上述情形中，有 $T$ 上层 的典范同构
\[
f^{*}(R^{r}\varphi_{*}F) \to R^{r}\varphi'_*(f'^{*}F).
\]
\end{theorem}

\begin{proof}
首先定义这个映射.由于 $f^{*}$ 与 $f_{*}$ 伴随，只需定义函子的态射
\[
R^{r}\varphi_{*} \to f_{*} \circ R^{r}\varphi'_* \circ f'^{*}.
\]
我们把它取为如下态射的复合：
\[
R^{r}\varphi_{*} \to R^{r}\varphi_{*} \circ f'_* \circ f'^{*}
\]
(由 $\mathrm{id} \to f'_* \circ f'^{*}$ 诱导),
\[
R^{r}\varphi_{*} \circ f'_* \circ f'^{*} \to R^{r}(\varphi \circ f')_* \circ f'^{*} = R^{r}(f \circ \varphi')_* \circ f'^{*}
\]
(映射 $R^{r}\varphi_{*} \circ f'_* \to R^{r}(\varphi \circ f')_*$ 由 Leray 谱序列给出)，以及
\[
R^{r}(f \circ \varphi')_* \circ f'^{*} \to f_{*} \circ R^{r}\varphi'_* \circ f'^{*}.
\]
于是，有 $T$ 上层 的典范映射 $f^{*}(R^{r}\varphi_{*}F) \to R^{r}\varphi'_*(f'^{*}F)$，因为它在茎上同构，它是同构.
\end{proof}

\begin{remark}\label{rem:17.4}
事实上，定理~\ref{thm:17.2} 是定理~\ref{thm:17.3} 中 $T \to S$ 取为 $s \to S$ 的特殊情形.
\end{remark}

\subsection*{几何中的紧合基变换定理}

回忆(AG 7.1) $k$ 上的代数簇 $X$ 称为完备的，如果对所有 $k$-簇 $T$，投影映射 $X \times T \to T$ 闭(即把 Zariski 闭集送到闭集).例如，$k$ 上的射影簇完备(AG 7.7)，而仿射簇完备，当且仅当它有限.

更一般地，簇(或概形)的正则映射 $X \to S$ 称为紧合的，如果对所有 $T \to S$，投影映射 $X \times_{S} T \to T$ 闭(且 $X \to S$ 分离).例如，有限态射，例如闭浸入，是紧合的(同上，6.24)，而 $k$-簇 $X$ 完备，当且仅当映射 $X \to P$(点)紧合.

回忆 $X_{\mathrm{et}}$ 上的层 $F$ 局部常值，如果存在覆盖 $(U_i \to X)_{i \in I}$，使得对所有 $i \in I$，$F|_{U_i}$ 常值.当 $F$ 茎有限时，这等价于存在单个有限满平展映射 $U \to X$，使得 $F|_{U}$ 常值.

遗憾的是，局部常值层的类对直像(甚至紧合映射的直像，甚至闭浸入的直像)不稳定，因此我们需要更大的类.

\begin{definition}\label{def:17.5}
$X_{\mathrm{et}}$ 上的层 $F$ 称为可构造的，如果
\begin{enumerate}
  \item 对每个闭浸入 $i \colon Z \hookrightarrow X$($Z$ 不可约)，存在非空开子集 $U \subset Z$，使得 $(i^{*}F)|_{U}$ 局部常值；
  \item $F$ 的茎有限.
\end{enumerate}
\end{definition}

\begin{remark}\label{rem:17.6}
可以证明每个挠层都是它的可构造子层的并.
\end{remark}

\begin{theorem}\label{thm:17.7}
设 $\varphi \colon X \to S$ 是紧合态射，$F$ 是 $X$ 上的可构造层.则对所有 $r \geq 0$，$R^{r}\varphi_{*}F$ 可构造，且对 $S$ 的每个几何点 $\bar{s} \to S$, $(R^{r}\varphi_{*}F)_{\bar{s}} = H^{r}(X_{\bar{s}}, F|_{X_{\bar{s}}})$.
\end{theorem}

这里 $X_{\bar{s}} \overset{\mathrm{def}}{=} X \times_{S} \bar{s}$ 是 $\varphi$ 在 $\bar{s}$ 上的几何纤维，$F|_{X_{\bar{s}}}$ 是 $F$ 在映射 $X_{\bar{s}} \to X$ 下的逆像.

\begin{corollary}\label{cor:17.8}
设 $X$ 是可分闭域 $k$ 上的完备簇，$F$ 是 $X$ 上的可构造层.
(a) 对所有 $r$，群 $H^{r}(X_{\mathrm{et}}, F)$ 有限.
(b) 对任何可分闭域 $k' \supset k$，$H^{r}(X, F) = H^{r}(X', F)$，其中 $X' = X_{k'}$.
\end{corollary}

\begin{proof}
(a) 考虑映射 $\varphi \colon X \to s$($X$ 到点).映射 $G \mapsto \Gamma(s, G)$ 把 $s_{\mathrm{et}}$ 上的可构造层范畴等同于有限 Abel 群范畴，并把 $R^{r}\varphi_{*}F$ 等同于 $H^{r}(X_{\mathrm{et}}, F)$.论断现在由定理推出，因为 $\varphi$ 紧合.

(b) 设 $\bar{s} = \operatorname{Spec} k$, $\bar{s}' = \operatorname{Spec} k'$.对应包含 $k \hookrightarrow k$(恒等映射)与 $k \hookrightarrow k'$ 的映射 $\bar{s} \to s$ 与 $\bar{s}' \to s$ 是 $s$ 的几何点.层在几何点 $\bar{t}$ 处的茎只依赖于 $t$，因此定理蕴含
\[
H^{r}(X, F) \xrightarrow{17.7} (R^{r}\varphi_{*}F)_{\bar{s}} \cong (R^{r}\varphi_{*}F)_{\bar{s}'} \xrightarrow{17.7} H^{r}(X_{k'}, F).
\]
\end{proof}

\begin{remark}\label{rem:17.9}
即使对 $\mathbb{Q}$ 上的点 $P$ 与层 $\mathbb{Z}/2\mathbb{Z}$，推论 (a) 的论断在去掉 $k$ 可分闭的条件后也假，因为在这种情况下 $H^{1}(P, \mathbb{Z}/2\mathbb{Z}) = H^{1}(P, \mu_2) \cong \mathbb{Q}^{\times}/\mathbb{Q}^{\times 2}$，它是由 $-1$ 与诸素数的类生成的 $\mathbb{F}_2$ 上的无限维向量空间.

推论 (b) 的论断在去掉 $X$ 完备的条件后也假，即使对 $\mathbb{A}^{1}$ 上的层 $\mathbb{Z}/p\mathbb{Z}$ 也假，因为 $\mathbb{Z}/p\mathbb{Z}$ 在 $\mathbb{A}^{1}_{k'}$ 上的逆像仍是 $\mathbb{Z}/p\mathbb{Z}$，且(见例~\ref{ex:7.9}b)
\[
H^{1}(\mathbb{A}^{1}, \mathbb{Z}/p\mathbb{Z}) = k[T]/(T^{p} - T) \neq k'[T]/(T^{p} - T) = H^{1}(\mathbb{A}^{1}_{k'}, \mathbb{Z}/p\mathbb{Z}).
\]
\end{remark}

\begin{theorem}\label{thm:17.10}
设 $\varphi \colon X \to S$ 紧合，对某个态射 $f \colon T \to S$，设 $X' = X \times_{S} T$：
\[
\begin{tikzcd}
X \arrow[d, "\varphi"'] & X' \arrow[l, "f'"'] \arrow[d, "\varphi'"] \\
S & T \arrow[l, "f"']
\end{tikzcd}
\]
对 $X$ 上的任何挠层 $F$，有典范同构
\[
f^{*}(R^{r}\varphi_{*}F) \to R^{r}\varphi'_*(f'^{*}F).
\]
\end{theorem}

\begin{proof}
这可以像定理~\ref{thm:17.3} 从定理~\ref{thm:17.2} 推出那样，从定理~\ref{thm:17.7} 推出：首先用函子的伴随性构造映射，然后用定理在茎上验证映射是同构.
\end{proof}

定理~\ref{thm:17.10} 就是紧合基变换定理.

\subsection*{关于证明的评注}

定理~\ref{thm:17.7} 与~\ref{thm:17.10} 是平展上同调基本定理中最难证明的.我解释为什么.本节中，我只考虑代数闭域上的簇.

设 $S$ 是簇，$S_{\mathbf{Et}}$ 是 $S$ 上的大平展景：$S_{\mathbf{Et}}$ 的底层范畴是 $S$ 上所有簇的范畴，覆盖族是 $S$-簇的平展映射的满族 $(T_i \to T)$.$S_{\mathbf{Et}}$ 上的层是满足所有覆盖的层条件(\S 5, S)的反变函子 $F \colon \mathrm{Cat}(S_{\mathbf{Et}}) \to \mathbf{Ab}$，即对每个正则映射 $T \to S$，$F$ 到 $T_{\mathrm{et}}$ 的限制都是层.

设 $f \colon S_{\mathbf{Et}} \to S_{\mathrm{et}}$ 是显然的连续态射.考虑 $S_{\mathbf{Et}}$ 上层 $F$ 的如下条件：
\begin{enumerate}
  \item $F \xrightarrow{\sim} f_{*}f^{*}F$；
  \item $f_{*}F$ 可构造.
\end{enumerate}
注意 $f_{*}F$ 是 $F$ 到 $S_{\mathrm{et}}$ 的限制.于是 (a) 说 $F$(以自然的方式)由它到 $S_{\mathrm{et}}$ 的限制决定，(b) 说 $F$ 到 $S_{\mathrm{et}}$ 的限制可构造.

\begin{proposition}\label{prop:17.11}
$S_{\mathbf{Et}}$ 上的层 $F$ 满足 (a) 与 (b)，当且仅当它由在 $S$ 上拟有限的簇 $F \to S$ 表示，即当且仅当对所有 $T \to S$，$F(T) = \operatorname{Hom}_S(T, F)$.
\end{proposition}

\begin{proof}
证明的讨论见 EC V.1.
\end{proof}

现在设 $\varphi \colon X \to S$ 是紧合映射，$F$ 是 $X_{\mathrm{et}}$ 上的可构造层.定义 $F^r$ 为 $S_{\mathbf{Et}}$ 上满足 $F^r|_{T} = R^{r}\varphi'_*(f'^{*}F)$(记号如定理中)的层.当 $T' \to T$ 不平展时，限制映射 $F^r(T) \to F^r(T')$ 由定理中的基变换映射给出.或者，可以定义 $F^r$ 为 $R^{r}\varphi_*(f'^{*}F)$，其中 $f'^{*}F$ 是 $F$ 在 $X_{\mathbf{Et}}$ 上的逆像，$R^{r}\varphi_*$ 在大平展景上计算.定理断言 $F^r$ 满足 (a) 与 (b).按命题，这等价于断言 $F^r$ 由在 $S$ 上拟有限的簇表示.

困难就在这里：我们有一个定义在 $S$ 上所有簇的范畴上的函子，希望证明它由在 $S$ 上拟有限的簇表示.这样的论断通常非常难证明.

这种情况证明的出发点是 Grothendieck 的如下定理：设 $\varphi \colon X \to S$ 是紧合正则映射；则 $R^{1}\varphi_{*}\mathbb{G}_{m}$ 在 $S_{\mathbf{Et}}$ 上由 Picard 概形 $\operatorname{Pic}_{X/S}$ 表示---它是 $S$ 上簇的无限并.由 Kummer 序列，发现
\[
R^{1}\varphi_{*}\mu_n = \operatorname{Ker}\left( \operatorname{Pic}_{X/S} \xrightarrow{n} \operatorname{Pic}_{X/S} \right),
\]
从而可表示.这个论证足以在 $\varphi$ 的纤维是曲线时证明定理，而一般情形对 $X$ 在 $S$ 上的相对维数作归纳证明.见 Artin 1973 的最后一章.
